\paragraph{Frage und Erwartung.} Kann das Ersatzmodell selbst bessere Entwürfe finden? Aus den sechs Gauß-Prozessen
wurde für jeden denkbaren Entwurf die kleinste Reserve vorhergesagt und davon eine Standardabweichung abgezogen
(vorsichtige Schätzung). Die vier Entwürfe mit dem höchsten Wert wurden voll gerechnet. Alle vier haben 24 Spulen,
einen Abstand Spule--Plasma von \SI{0,21}{m}, kurze Spulen (\SIrange{3,9}{4,1}{m} je Spule) und die dickste
Außenwand; sie unterscheiden sich vor allem in Zahl und Steifigkeit der Verbinder. Die vorsichtige Vorhersage lag bei
0,04--0,07. Erwartet war (Ereignisprotokoll F-142): 1--3 der 4 zulässig, Spannung $\pm30$~\% neben der Vorhersage,
keiner deutlich robuster als 0169.

\begin{table}[H]\centering\scriptsize\begin{tabular}{@{}lrrrrrrrrl@{}}\toprule
Vorschlag & $E_\text{V}$ & V./Paar & Masse & $\sigma_{99}$ & $\sigma_{99}$ fein & Dehnung & $B\cdot n$ & \multicolumn{2}{c}{Robustheit} \\
 & GPa & & kt & MPa & MPa & \% & ‰ & vorherges. & gerechnet (Engpass) \\\midrule
V0 & \num{55} & 13 & \num{37.8} & \num{306} & \num{385} & \num{0.16} & \num{1.52} & 0,07 & 0,23 (Platz Package) \\
V1 & \num{77} & 13 & \num{33.1} & \num{220} & \num{366} & \num{0.12} & \num{1.87} & 0,06 & 0,31 (Platz Package) \\
V2 & \num{81} & 9 & \num{26.4} & \num{851} & -- & \num{0.54} & \num{2.19} & 0,06 & -0,79 (Dehnung Wickelpaket), unzulässig \\
V3 & \num{176} & 11 & \num{29.1} & \num{272} & \num{394} & \num{0.13} & \num{1.96} & 0,04 & 0,35 (Feldfehler B·n) \\
\bottomrule\end{tabular}
\caption{Vorschläge des Ersatzmodells.
$E_\text{V}$ Verbindersteifigkeit, V./Paar Verbinder je Spulenpaar, „fein“ = 1,5-mal mehr Stationen.
Robustheit = kleinste Reserve über alle acht Grenzen; „vorherges.“ = Vorhersage minus eine Standardabweichung.}
\end{table}

\paragraph{Ergebnis.} Drei der vier Vorschläge sind zulässig, die gerechnete Robustheit (0,23--0,35) liegt sogar
weit über der vorsichtigen Vorhersage. V3 (11 steife Verbinder je Paar, \SI{176}{GPa}) erreicht 0,35 und ist im
Standardnetz der robusteste aller bisher gerechneten Entwürfe; Engpass ist dort der Feldfehler, nicht die Spannung.
V2 dagegen, mit nur 9 Verbindern je Paar, scheitert deutlich: $\sigma_{99}$ \SI{851}{MPa}, Dehnung des Wickelpakets
0,54~\%, Verformung \SI{177}{mm}. Das Ersatzmodell hat diesen Ausreißer nicht kommen sehen.

\paragraph{Netzprüfung.} Wie bei den doe1-Kandidaten steigt $\sigma_{99}$ mit 1,5-mal mehr Stationen deutlich: V3
\SI{272}{} $\to$ \SI{394}{MPa} ($+45$~\%), V1 \SI{220}{} $\to$ \SI{366}{MPa} ($+66$~\%), V0 \SI{306}{} $\to$
\SI{385}{MPa} ($+26$~\%). Alle bleiben zulässig, doch V3 fällt damit auf eine Robustheit von 0,21, weil jetzt der
Stahl begrenzt. 0169 (\SI{319}{MPa} im feinen Netz, Robustheit 0,29) bleibt der robusteste Entwurf.

\paragraph{Deutung.} Das Ersatzmodell ist für die Vorauswahl brauchbar: Es findet zulässige Gebiete zuverlässig
(Package, Kernheizung und Feldfehler sagt es gut vorher). Für die Spannung reicht es nicht; ein einzelner Vorschlag
kann weit danebenliegen, und die Netzabhängigkeit der Spitzen sieht es gar nicht. Auffällig ist, dass 0169 der
einzige Entwurf bleibt, dessen Spannung beim feineren Netz kaum steigt. Bei 0169 verschmelzen die Verbinder zu fast
durchgehenden Platten (Abschnitt~\ref{sec:3d}). Eine Vermutung, noch ungeprüft: Weniger freie Verbinderenden
bedeuten weniger Steifigkeitssprünge und damit Spitzen, die mit dem Netz nicht wachsen.
